\documentclass[10pt,a4paper]{amsart}
\usepackage[margin=19mm]{geometry}
\usepackage{amsmath,amssymb,mathtools}
\usepackage[scr=rsfs,cal=euler]{mathalfa}
\usepackage{xfrac}
\usepackage[unicode,hidelinks,bookmarks=false]{hyperref}
\newcommand{\ie}{i.e.\ }
\newcommand{\cf}{cf.\ }
\newcommand{\resp}{respectively\ }
\newcommand{\Subsection}{\subsection}
\providecommand{\nobreakdash}{\nobreak\hbox{-}}
\setlength{\emergencystretch}{2em}
\multlinegap0pt
\usepackage{bm}
\usepackage{bbm}
\usepackage{dsfont}
\usepackage{xspace}
\usepackage{cancel}
\usepackage{mathtools}
\usepackage{amsthm}
\makeatletter
\@ifundefined{theorem}{\newtheorem{theorem}{Theorem}}{}
\@ifundefined{lemma}{\newtheorem{lemma}[theorem]{Lemma}}{}
\@ifundefined{proposition}{\newtheorem{proposition}[theorem]{Proposition}}{}
\makeatother
\title[Central polynomials of minimal degree for matrices]{Central polynomials of minimal degree for matrices}
\author{Vesselin Drensky, Boyan Kostadinov}
\date{Source: arXiv:2601.07750v2 (CC BY 4.0)}
\begin{document}
\maketitle

\noindent\emph{Source and licence.} Central polynomials of minimal degree for matrices. Version arXiv:2601.07750v2 (CC BY 4.0), distributed under the Creative Commons Attribution 4.0 International licence, as recorded in the arXiv licence metadata for this exact version and shown on the abstract page https://arxiv.org/abs/2601.07750v2. Full rights notice: licenses/arxiv-2601.07750-CC-BY-4.0.txt.\par\smallskip

\noindent\textit{Editorial scope.} Counted unit: the Proposition whose source label is \texttt{no central polynomials of degree 12} in the article (source0.tex lines 916--918, source label \texttt{no central polynomials of degree 12}), delivered together with the article's own complete original proof of it (source0.tex lines 920--995, one \texttt{proof} environment of 76 lines). No mathematics is written, repaired or re-derived: statement and proof are the article's own text line for line. Required context retained uncounted: no central polynomials of degree 10 (lines 735--737); HWV of degree 12 (lines 862--914). Every definition, display and label of those lines is retained unchanged. Omitted from this package and not redistributed: 1--734, 738--861, 915--915, 919--919, 996--1283. Nothing inside a retained excerpt is omitted or altered: every retained body is delivered exactly as the article's lines are written, so no omission receipt of any kind accompanies this package. Citations: the retained excerpts carry no citation command at all, so no bibliography entry of the article is transcribed and no reference is redistributed. Reference adaptation: none was needed and none was made; every reference inside the delivered unit resolves inside it, so no unresolved reference or citation is produced. Printed slips kept literal: no printed word, symbol, hypothesis or number of the counted statement or of its proof was altered. Layout and class: the article's own class is \texttt{amsart} with packages including amscd, amssymb; the wrapper is the lane's pinned \texttt{amsart} editorial wrapper at 10pt on A4, which restates the theorem-like environments the retained text needs and adds the article's own macro definitions that the retained text needs (none), with extra wrapper packages (bm, bbm, dsfont, xspace, cancel, mathtools). The delivered numbering is the unit's own. Assets: no retained span contains a figure environment, a graphics-inclusion command, a PDF-page inclusion, a table or a TikZ picture; no image file is copied into this package, the package declares no asset dependency and no image file is redistributed with it. Licence: version arXiv:2601.07750v2 of this article is distributed under the Creative Commons Attribution 4.0 International licence, as recorded in the arXiv licence metadata for this exact version and shown on the abstract page https://arxiv.org/abs/2601.07750v2. A full rights notice is delivered at licenses/arxiv-2601.07750-CC-BY-4.0.txt. Characters: every retained excerpt is pure ASCII, so the delivered engine, XeLaTeX with its default Latin Modern text font, reports no missing character.
\begin{proposition}\label{no central polynomials of degree 10}
The algebra $M_4(K)$ does not have central polynomials $c(x,y)$ of degree $10$.
\end{proposition}

\begin{lemma}\label{HWV of degree 12}
{\rm (i)} The following polynomials are highest weight vectors in $\text{\rm Com}_{m_1}\cdots \text{\rm Com}_{m_k}$, $m=(m_1,\ldots,m_k)$, for $\lambda=(7,5)$:
\[
(5,3,2^2):w_1(x,y)=([y,x,x,x,x][y,x,y]-[y,x,x,x,y][y,x,x])[y,x]^2, %\in\text{\rm Com}_5\text{\rm Com}_3\text{\rm Com}_2^2,
\]
\[
(4^2,2^2):w_2(x,y)=([y,x,x,x][y,x,x,y]-[y,x,x,y][y,x,x,x])[y,x]^, %\in\text{\rm Com}_4^2\text{\rm Com}_2^2,
\]
\[
(4,3^2,2):w_3(x,y)=([y,x,x,x][y,x,y]-[y,x,x,y][y,x,x])[y,x,x][y,x], %\in\text{\rm Com}_4\text{\rm Com}_3^2\text{\rm Com}_2,
\]
\[
(4,3^2,2):w_4(x,y)=[y,x,x,x]([y,x,x][y,x,y]-[y,x,y][y,x,x])[y,x], %\in\text{\rm Com}_4\text{\rm Com}_3^2\text{\rm Com}_2,
\]
\[
(3^4):w_5(x,y)=([y,x,x][y,x,y]-[y,x,y][y,x,x])[y,x,x][y,x,x],%\in\text{\rm Com}_3^4,
\]
\[
(4,2^4):w_6(x,y)=[y,x,x,x][y,x]^4,%\in\text{\rm Com}_4\text{\rm Com}_2^4,
\]
\[
(3^2,2^3):w_7(x,y)=[y,x,x]^2[y,x]^3. %\in\text{\rm Com}_3^2\text{\rm Com}_2^3.
\]

{\rm (ii)} The following are highest weight vectors for $\lambda=(6,6)$:
\[
(4^2,2^2):w_1(x,y)=([y,x,x,x][y,x,y,y]-[y,x,y,x][y,x,x,y]-
\]
\[
-[y,x,x,y][y,x,y,x]+[y,x,y,y][y,x,x,x])[y,x]^2,
\]
\[
(4,3^2,2):w_2(x,y)=([y,x,x,x][y,x,y][y,x,y]-[y,x,y,x][y,x,x][y,x,y]-
\]
\[
-[y,x,x,y][y,x,y][y,x,x]+[y,x,y,y][y,x,x][y,x,x])[y,x],
\]
\[
(3^4):w_3(x,y)=([y,x,x][y,x,y]-[y,x,y][y,x,x])^2,
\]
\[
w_4(x,y)=[y,x,x][y,x,x][y,x,y][y,x,y]-[y,x,y][y,x,x][y,x,x][y,x,y]-
\]
\[
-[y,x,x][y,x,y][y,x,y][y,x,x]+[y,x,y][y,x,y][y,x,x][y,x,x],
\]
\[
(3^2,2^3):w_5(x,y)=([y,x,x][y,x,y]-[y,x,y][y,x,x])[y,x]^3,
\]
\[
(2^6):w_6(x,y)=[y,x]^6.
\]
\end{lemma}

\begin{proposition}\label{no central polynomials of degree 12}
The algebra $M_4(K)$ does not have central polynomials and polynomial identities of degree $12$ in $K\langle x,y\rangle$.
\end{proposition}

\begin{proof}
Again, we follow the main steps of the proof of Proposition \ref{no central polynomials of degree 10}.
It is sufficient to consider the case $\lambda=(7,5)$ and $\lambda=(6,6)$.

(i) The case $\lambda=(7,5)$. For
\[
(m_1,\ldots,m_k)=(4^2,2^2),(3^2,2^3),(2^6)
\]
there is only one $GL_2(K)$-submodule $W(7,5)$ in the products of commutators
\[
\text{\rm Com}_{m_1}\cdots\text{\rm Com}_{m_k},m_1+\cdots+m_k=12,k\geq 4.
\]
The number of direct summands $W(7,5)$ is equal to 2 in the cases
\[
(m_1,\ldots,m_k)=(4,3^2,2),(3^4).
\]
The corresponding highest weight vectors are given in Lemma \ref{HWV of degree 12} (i).
Depending on the positions of the commutators of length $m_1,\ldots,m_k$ the number of linearly independent highest weight vectors is:
\begin{itemize}
\item $(5,3,2^2)$, $w_1(x,y)$: 12,
\item $(4^2,2^2)$, $w_2(x,y)$: 6,
\item $(4,3^2,2)$, $w_3(x,y),w_4(x,y)$: 24,
\item $(3^4)$, $w_5(x,y),w_4(x,y)$: 3,
\item $(4,2^5)$, $w_6(x,y)$: 5,
\item $(3^2,2^3)$, $w_7(x,y)$: 10.
\end{itemize}
Totally we have 60 linearly independent highest weight vectors.
We consider
\[
w_{(7,5)}(x,y)=\sum_{i=1}^{60}\xi_iw^{(i)}(x,y)
\]
with unknown coefficients $\xi_i\in K$, $i=1,\ldots,60$. As in Proposition \ref{no central polynomials of degree 10}
we replace $x$ by the diagonal generic matrix $u$. Now we replace $y$ by the matrix
\[
v_2=\left(\begin{matrix}1&1&v_{13}&1\\
0&1&1&1\\
1&1&1&1\\
1&1&0&0\\
\end{matrix}\right).
\]
We consider the entry $z_{12}$ of
\[
w_{(7,5)}(x,y)=\sum_{a,b=1}^4z_{ab}E_{ab}.
\]
It is in the form
\[
z_{12}=\sum_{i,j}f_{ij}(\xi_1,\ldots,\xi_{60})u_1^{i_1}u_2^{i_2}u_3^{i_3}u_4^{i_4}v_{13}^j.
\]
We have computed that the homogeneous linear system consisting of all
\[
f_{ij}(\xi_1,\ldots,\xi_{60})=0
\]
has only the zero solution and this shows that $M_4(K)$ does not have polynomial identities and central polynomials for $\lambda=(7,5)$.

(ii) The case $\lambda=(6,6)$. As in the case of the proof of (i), the corresponding highest weight vectors are given in Lemma \ref{HWV of degree 12} (ii).
Depending on the positions of the commutators of length $m_1,\ldots,m_k$ in the product of $\text{\rm Com}_{m_1},\ldots,\text{\rm Com}_{m_k}$,
the number of linearly independent highest weight vectors is:
\begin{itemize}
\item $(4^2,2^2)$, $w_1(x,y)$: 6,
\item $(4,3^2,2)$, $w_2(x,y)$: 12,
\item $(3^4)$, $w_3(x,y),w_4(x,y)$: 2,
\item $(3^2,2^3)$, $w_5(x,y)$: 10,
\item $(2^5)$: $w_6(x,y)$: 1.
\end{itemize}
Totally we have 31 linearly independent highest weight vectors. As in (i), we consider the linear combination $w_{(6,6)}(x,y)$ of these 31 polynomials,
evaluate it for $x=u$, $y=v_2$ and solve the homogeneous linear system obtained from the entry $z_{12}$. The system has only the trivial solution which shows
that $M_4(K)$ does not have polynomial identities and central polynomials for $\lambda=(6,6)$.
The same result can be obtained if we use the entry $z_{12}$ of the evaluation $w_{(6,6)}(u,v_3)$ instead of the evaluation $w_{(6,6)}(u,v_2)$, where
\[
v_3=\left(\begin{matrix}0&1&v_{13}&0\\
0&1&1&0\\
0&1&0&1\\
1&0&0&0\\
\end{matrix}\right).
\]
\end{proof}

\end{document}
